\chapter{如何讓人妻愛上你}

\section{泡妞是一門科學}

其實量化交易的核心就一句話：交易是一門科學。只要你相信科學，相信市場是可以（從概率上）預測的，然後用科學的實證的方法研究，你就懂量化了。無論你是做高頻交易還是低頻交易，無論你是用統計模型還是AI，無論你是用自動下單系統還是手動把房子掛出去賣掉都不重要，科學的交易者必然把信「纏論」或者「一勞永逸地把大A打到10000點」的前現代土鱉踩在泥土裏揉成渣。

我每次用科學的方法論在市場上輕鬆賺錢的時候，感覺就像是那些拿着火槍，身着鐵甲，以一敵百征服、屠殺印加帝國的那168名西班牙征服者。而歷史上對印加帝國的征服，又何嘗不是掌握了科學的一方對前現代的一方單方面的碾壓與屠殺呢？

我看今天絕大多數中國人，接受了中學甚至大學教育，但是卻從來沒有真正掌握科學精神和科學方法論。（甚至很多反而是被學校灌輸了很多錯誤觀念。）一遇到教科書裏沒有出現過的問題，就求助於迷信和僞科學。

我的大姨，中山大學畢業的金中前校長李麗麗在2007年的時候拿了我媽幾十萬，重倉中石油血本無歸。她的投資決策是基於某種信仰而不是科學。而我在的一個上海中學校友羣裏面，一個上海中學和上海交通大學畢業的上海妹妹，每天對着手機研究我的星座、命盤和八字。這種愚蠢的行爲，讓人感慨她受到的科學教育都扔黃浦江裏面去了。

\textbf{要學習怎麼泡妞，第一步，就是要認識到，泡妞是一門科學，女人的心理和行爲是可預測的，甚至是非常容易預測的。}

你需要把社會尤其是女人給你灌輸的「女人是神」這個信仰或者說思想鋼印給去掉。前面說過，這個信仰甚至是源於西方的天主教，根本不是中國文化。但是現在，有一羣人在拼盡全力維護「女人是神」這個概念，一旦你試圖開始審視、研究女人，你就像《楚門的世界》（The Truman Show）裏面的楚門一樣，會受到周圍所有人的攻擊：會說你是不尊重女性、物化女性，等等等等。但是你如果信了這些女人和龜男的話，你就會像A股市場裏面信「10000點」的韭菜們一樣，血本無歸，甚至爆倉跳樓。

\section{泡妞的科學原理}

這一節可能會涉及一些比較專業的知識，主要集中在博弈論（game theory）。我會盡可能保留專業術語（以便有興趣的讀者深入研究）的同時，用通俗易懂的語言解釋，讓非專業的普通讀者也能理解理論。只有理解了科學的理論，才能夠舉一反三，在實操中隨機應變，立於不敗之地。

關於泡妞的科學的基礎，可以總結爲以下四個公理。

\subsection{李新野泡妞第一公理：感覺人公理}

\begin{axiom}\label{a1}
在自由狀態下，女人絕大多數決策都是基於感覺，而感覺是基於本能，而本能是基於人類上百萬年原始社會中兩性博弈進化出來的納什均衡策略\footnote{如果你沒學過納什均衡，可以把納什均衡策略理解爲（局部穩定）最優策略。}。
\end{axiom}

現代進化論告訴我們，包括人類在內的動物，在進化論的篩選下，一切的行爲，都需要服務於生存與繁衍。但是這顯然無法解釋很多問題。比如爲什麼現在中國美國的沿海大城市集中了那麼多單身女性，寧願斷子絕孫也不嫁人。這就回到了一個其實很容易觀測到的事實：女人絕大部分的決策，都是基於自身的感覺。

那如何連接上這兩個看似矛盾的理論呢？其實很簡單，就是女人的感覺，其實就是上百萬年中進化出來的最有利於女人生存繁衍的本能。最近5000年的父系傳統文明的科技發展實在是太迅猛了，而人類的本能還未進化到適應現代高科技社會。而女人在失去了三從四德的規訓之後，靠着感覺，重新跟隨原始社會的最優策略在生活。

如果說，現代經濟學是基於用博弈論分析「理性經濟人」（economic man）\footnote{「理性經濟人」也翻譯爲「理性人」或者「經濟人」。指的是在經濟學研究中，我們假設所有的人都是自私、無情感地追求自身經濟利益最大化。雖然現實中金錢並不是人類唯一的行爲動機，但是這種簡單的假設可以簡化數學模型，讓博弈論分析變得可行且簡單清晰。}。那麼我們可以把現代社會的女人稱之爲「感覺人」。我們完全可以直接假設她們做的所有事情，都是爲了讓她們感覺更好。而把滿足感覺（也就是原始社會的最佳策略）作爲「效用」（utility）\footnote{「效用」是指博弈論裏每個參與者（玩家，player）試圖最大化的一個數值。在理性經濟人假設下，效用等同於金錢。而在感覺人假設下，效用等同於感覺，也就是對原始社會最佳男女博弈策略的接近程度。}，便可以開始用博弈論研究感覺人（女人）的行爲了。

\subsection{李新野泡妞第二公理：一夫多妻公理}

\begin{axiom}\label{a2}
一夫多妻是人類歷史的常態，女人慕強和向上擇偶是納什均衡策略。
\end{axiom}

首先，結合公理\ref{a1}，很容易也推導出以下定理。

\begin{theorem}\label{t1}
女人慕強和向上擇偶是天性。
\end{theorem}

這裏說一夫多妻\footnote{指一個男人有多個女配偶，包含一夫一妻多妾。}是「常態」，並不是講所有男人都有多個配偶，根據費希爾原理（Fisher's principle）\footnote{指根據自然選擇理論，包括人類在內雌雄異型不大的物種，性別比會接近於1:1。否則數量少的性別有進化優勢，進而恢復平衡。}，人類性別比接近1:1，每個男人有多個配偶在數學上是不可能的。

這裏的常態也不是說每個女人都要做妾。實際上，我翻看各種資料，無論是中國古代的史料，還是現在穆斯林國家和撒哈拉以南非洲的社會狀況，即使一夫多妻完全合法，也僅僅有30\%左右的男人會有兩個或以上的配偶。

但是，只要社會中常態存在一夫多妻的家庭，只要進入一夫多妻關係是女人的選擇之一，那麼這就會深刻改變女人的感覺和擇偶策略。

我們做一個思想實驗，如果人類有這麼一個開關，如果和第二個異性發生性關係，就會立刻死掉。或者更溫和點，人類只能和第一個發生性關係的人產生後代。那麼經過千萬年的進化，無論是男人還是女人，都會傾向於找一個優秀程度在同性中的百分位和自己差不多的異性，安心結合繁衍後代。但是這明顯不符合現實。

而因爲超過一百萬年來，人類一夫多妻是常態，女性主要負責撫養後代。那麼男性最佳的擇偶策略自然是會傾向於多找幾個配偶，即使配偶優秀程度差於自己，但是能增加後代數量。而女性，因爲一生生育數量有限，且需要爲撫養後代付出更多，所以女性的最佳擇偶策略便是把有限的生育資源留給更優秀的男性。盡量生下與優秀男性的共同後代。

所以這一切就解釋得通了，因爲一夫多妻是常態，所以在上百萬年的進化下，男女的天性符合一夫多妻的納什均衡策略：男人只要看到年輕漂亮忠誠的女人，內心都不會拒絕。但是要讓一個女人跟一個沒有她優秀的男人結合甚至繁衍後代，女人的內心是會極其痛苦的，覺得自己在這個一夫多妻的擇偶博弈中失敗了，基因被污染了，沒法獲得優秀的後代。

結合公理\ref{a1}和公理\ref{a2}，可以得出以下推論。

\begin{corollary}\label{c1}
女人慕強，道德上等價於男人多偶。
\end{corollary}

女性慕強和男性多偶，本質上都是在現代社會釋放原始社會的天性。但是現在很多華師妹竟然可以心安理得地覺得自己和清華男是配得上的，然後社會都去合理化這種慕強擇偶。但是一個清華男說要有找小妾或者「細姨」\footnote{閩南語的小老婆}，全社會都會去譴責，這是很匪夷所思的事情。

然後結合公理\ref{a1}和公理\ref{a2}，還能得到下面這個很讓人震驚的定理。

\begin{theorem}\label{t2}
任何同時解放婦女和實行一夫一妻制的民族，必然消亡。
\end{theorem}

在上萬年的一夫多妻下，女人進化出了慕強的本能。所以女人和（在擇偶吸引力上）比自己弱甚至與自己相當的男性交配、繁衍這件事，是發自內心的厭惡和痛苦。但是因爲一夫一妻制，她們看得上的男人，根本看不上她們，不可能和她們進入一夫一妻的關係。所以女人的決策基於自身感覺，爲了避免這種下嫁或者平嫁的痛苦，女人便會寧願選擇不婚不育，也不與看不上的男性結婚、交配、生育。

所以解放婦女加上一夫一妻制到最後，會有大約一半的女人終身不婚不育。這個民族自然是要消亡的。

現在中國的結婚率，還是因爲有長輩催婚撐上來的，但是總和生育率也小於1了，中國人會以指數速度絕種。而美國，基本沒有保進步守的人，所以稍微好點。一方面，美國有部分人有守可保，可以退回原教旨基督教、正統猶太教的傳統，不像中國，孔子和儒家和三從四德已經被完全打倒了。另一方面，美國在自由主義上，又正在慢慢進步到邏輯自洽的全面解放，所以非婚生子幾乎完全不受歧視。也算是事實上部分開放了一夫多妻。現在美國大約一半的嬰兒是非婚生子——這還不包括婚內女人和其他男人生下的小孩。靠這兩點，美國現在的總和生育率顯著超過中國，但是還是遠小於更替生育率。

所以，如果你用科學的方法來看，無論是中國、美國、北歐還是日本，在解決少子化的問題上的路徑都走錯了。這些國家都試圖用提高女性權利和福利，也就是賄賂女性的方式來提高生育率。但這注定是徒勞無功的，因爲女性對平嫁和低嫁的厭惡和痛苦，遠遠超過國家能給的幾千幾萬生育補貼能給女性帶來的快樂。所以凌菲菲和許潔被迫和她們看不上的李松堅上牀，只能吸乾李松堅幾億幾十億財產報仇，甚至故意把我的姐姐弄殘廢。即使這樣，她們也無法真正解氣。\footnote{見附錄收錄文章\hyperref[3faces]{《我的父親李松堅》}和\hyperref[sister]{《我的姐姐李新瑩》}。}

所以，要提高生育率，只能一方面堵，也就是恢復傳統道德，也就是按頭教化逼迫女人必須爲對她有恩的男人生育；另一方面疏，就是恢復一夫多妻制，讓一部分女人慕強心理比較強而嫉妒心比較弱的女性，可以有做小妾的選擇。

\subsection{李新野泡妞第三公理：女人無愛情公理}

\begin{axiom}\label{a3}
女人只對子女有愛，對男人只有爲獲取基因的迷戀和獲取供養和保護的索取。
\end{axiom}

這裏「愛」定義爲無條件付出的感情。而如果把「愛情」定義爲男女之間無條件付出的感情的話，那麼很容易得到以下推論。

\begin{corollary}\label{c2}
女人沒有愛情。
\end{corollary}

根據公理\ref{a1}，女人的天性和感覺，其實是服務於原始社會的生存和繁衍的。所以很好理解女人只會對自己的子女有無條件付出的感情。

而女人對男人的感情，其實服務於女人對男人的兩大需求。第一個需求是提供優質基因。所以女人在看到本能認爲的基因優秀的男人的時候，會陷入一種瘋狂的迷戀，想要立刻和對方交配、繁殖。

但是這種迷戀並不是愛情。女人只會想和你上牀，並不一定想和你發展長期關係，也很難給你真正的錢。\footnote{如果你想要從女人身上拿錢，你需要先讓女人迷戀，再讓女人得不到你，這時候女人就可能打錢了。這就是牛郎店賺錢的套路。但是你一旦跟她上牀了，她就不可能給你錢了。}而且，這種迷戀來得快，走得也快。也就是女人口中所謂的「上頭」和「下頭」。畢竟已經上牀了，女人想要的種已經得到了，沒有必要繼續在這個男人身上浪費時間了。

女人對男人的第二個需求，就是要求男人提供供養和保護的需求。當然根據公理\ref{a1}，女人這個需求也是有對應的感覺的，而這個感覺就是女人口中的「安全感」。

很多男人不理解，女人其實是被詛咒的，她們無時無刻不處於缺乏安全感的焦慮之中。這種焦慮比男人缺乏性生活還要痛苦。之前人人網（校內網）網紅李碩曾經將這種焦慮感稱之爲「陰莖缺失焦慮綜合症」，也算是非常精闢而準確的稱謂了。女人是永遠無法獨自承擔自身命運的，也無法獨立面對這個世界。女人需要不斷地確認是否能綁定一個強大的男人，讓他有意願源源不斷地提供供養和保護。

其實除了凌菲菲等極少數聰明的女人，絕大部分女人對男人的控制和吸血，主要並不是精確計算的結果，而是被「安全感缺失」痛苦所控制。也就是一旦男人沒有順從女人，沒有給錢，沒有儀式感，女人就會痛苦，沒有安全感，然後開始作鬧。而當你通過語言、行動、儀式感等方式緩解了女人的痛苦，女人便會停止作鬧，甚至會給你性交等好處。

這兩種感覺的對象，有人稱之爲「情人」和「供養者」，或者「基因提供者」和「供養者」。那麼這兩者，到底哪種是愛情呢？只能說見仁見智了。前者是瘋狂的迷戀，後者是緩解了女人的焦慮之後，產生的依戀。但是其實說實話，這兩者本質上都不是愛情。畢竟前者本質上只是要基因，後者本質上只是想要供養與庇護，最多在滿足之後給予身體接觸。都不是像基督教要求的教徒對上帝的那種無條件的愛與崇拜。

而男人其實是和女人不一樣的。就是男人其實是有一種無條件利女的蜂巢思維（hive mind）。也就是說，因爲女人是族羣繁育下一代的主體，所以保護女人，即使不是自己的配偶的女人，也對幫助族羣整體的繁衍有利。就像蜂羣裏面，放棄了個體繁衍的工蜂和雄蜂，犧牲自己，爲了蜂后的繁衍付出一切。其實男人和蜂羣裏面的工蜂和雄蜂一樣，有爲女人無條件犧牲的本能。這就是爲什麼會有泰坦尼克號一等艙男乘客生存率比三等艙女乘客還低——一個男人，即使努力到了買得起一等艙，在社會看來，生命的價值不如一個三等艙的底層女人。這種人類整體無意識地對男人生命的漠視，被稱之爲男性可棄置性（male disposability），而人類這種對女人的超額關注，被稱爲女本位主義（gynocentrism）。

而男人這種無條件犧牲自己造福女人的利女本能，也就是男性對女性的愛情。而女性的安全感缺失，裝可愛，裝可憐和作鬧之所以有用，也就是靠觸發身邊男性的這種利女本能（愛情），完成資源的掠奪和轉移。

因爲女人無論是對前者迷戀，還是對後者的依戀，出發點都是自身當下的情緒。也就是說，即使身邊的男人不是女人子女的父親，女人也會繼續釋放不安全感，作鬧，掠奪身邊男性的資源。這個當下的情緒，讓女人無論子女的血緣關係掠奪每個身邊的男人，這也是在原始社會下博弈的最優解。又印證了公理\ref{a1}。

比如最近有很多人妻出軌小孩非親生的案件，我很喜歡看這裏面對人妻的採訪，她們的回答其實非常有意思。有一個人妻說：「我就是出去放縱了一次，有必要這麼生氣嗎？」還有一個人說：「我都跟你那麼多年了，他都叫你爸爸那麼多年了，你就這麼拋棄我們你有沒有良心。」很多人看到這些言論非常氣憤，但是你看到這裏，你就會覺得這些女人其實誠實得甚至有點可愛。對那個情人，那個女人的迷戀的感情來得快，去得也快，也成功地獲得了他的基因。而對於她的丈夫，她還是把他當作長期緩解焦慮感和不安全感的對象，確實是把他當依靠對象和丈夫看的。所以她本能地對丈夫離開感到憤怒和不滿。

\subsection{李新野泡妞第四公理：間接探測公理}

\begin{axiom}\label{a4}
女人對男人的迷戀感覺基於對男人（在原始社會）的價值的間接判斷。
\end{axiom}

女人的直覺，其實是進化成一個「高價值男人探測器」或者說是「優質男人探測器」的。但是這種探測並不是直接評估價值，而是通過間接的線索去判斷這個男人的價值。

過去兩年，我經常拿着某個有100個比特幣的錢包給女人們看，作爲一種社會實驗。那些女人們，就算讀到985大學，看到那串餘額數字和數字背後的雜湊值，都是面無表情。那些比較有禮貌的或者比較有知識的，會「哦」一聲回應我。而那些既沒禮貌也沒知識的呢，就會露出一種對油膩中年暴發戶「嫖客」的那種厭惡表情。

而與之形成鮮明對比的是，有一天晚上，我在上海和一個江西九江的妹妹約會，我們在盧灣區的日月光中心喫了一頓美味的日式烤肉之後，帶她到上海西郊的一棟獨棟別墅約會。在這種浪漫氛圍的烘托之中，一切的進展都是那麼簡單而順理成章。

但是這種浪漫氛圍其實花不了多少錢！靠着我錢包裏那100個比特幣，即使不算未來比特幣的增值和上海房價的持續崩盤，我也能夠把自己複製四份，分別約一個彭澤妹妹、一個潯陽妹妹、一個柴桑妹妹和一個永修妹妹，然後分別到盧灣日月光、徐匯日月光、陸家嘴國金中心商場和靜安嘉里中心，再分別喫一頓日式烤肉、一頓韓國烤肉、一份美式烤牛排、一份巴西烤牛排，再分別帶她們到浦東、浦西、蘇（州河）南、蘇北四套獨棟別墅裏面約會上牀過夜。我可以這樣持續三十年，直到我完全陽痿爲止，然後這100個比特幣還有剩！

但是女人就是這樣子，她們不會爲100個比特幣對我產生迷戀，但是迷醉的燭光晚餐和燈火輝煌的大別墅，可以讓她們陰道溼潤並開始排卵。即使前者的價值是後者的數萬倍。

但是某種意義上，如果女性進化出對100個比特幣的發情本能，那20年前女人是不是會因爲沒有人有比特幣而無法發情而絕種？20年後的女人，是不是會因爲全世界除了20年後的我之外幾乎沒人有100個比特幣而絕種？所以女人這種「廉價七成正確」的擇偶本能，其實也就可以理解了。至於剩下的三成，屬於必要的基因缺陷消耗品。

當然女人對男人價值的間接探測，不只是燭光晚餐和大別墅那麼簡單。在後面的「吸引力法則」一節中，會介紹女人各種各樣的「價值探測器」並教大家如何黑進（hack）女人這套上萬年沒有更新過的擇偶軟件。

\section{AI時代的母猴子}

回想我自己的人生，我走的很多彎路，感受到的很多痛苦，都是因爲錯誤地把現代女人當作和我一樣的君子，或者我母親一樣的賢惠的女人。

在清華和MIT，我師從姚期智教授和Silvio Micali教授兩位圖靈獎得主學習、研究博弈論\footnote{筆者在清華和MIT的研究方向是機制設計（mechanism design），一個理論電腦和經濟學、博弈論的交叉學科，旨在研究如何設計商業、政治等博弈機制以優化利潤（profit）、社會福祉（social welfare）等不同效用。}。之後我再用我圖靈獎水平的大腦，和我的專業知識回頭審視現代女人，構建了這套理論，也就看開了：其實現代女性，不過是一羣活在AI時代的雌性靈長類動物（母猴子）罷了。

在太平洋兩岸的最頂級的科技和AI公司裏面，漢人男性貢獻了大約3/4的開發工作。而華人女性，卻每天打拳，用最惡毒的語言詛咒漢人男性。但是當你看完之前的理論分析，用原始人的眼光去看那羣國女，也就釋然了。如果你用邏輯去分析她們的語言，自然是狗屁不通。但是你如果用「感覺」去分析，一切就豁然開朗了。女權分子說的所有話，其實都是要釋放女性兩個底層的情緒：厭惡普通男性、安全感缺失。

前者的典型代表就是鄉下女拳頭子楊笠的「男人那麼普通又那麼自信」和官方女拳頭子張桂梅的「溪流」、「溝壑」、「草芥」、「懦夫」。網絡上現代女性鋪天蓋地的對催婚的反感也是一樣。上百萬年一夫多妻制的篩選下，讓女人們忍不住幻想一個清華畢業、美國歸國、白手起家、身家百億、然後還敢指着她們的鼻子罵她們是母猴子的強者從天而降到虹橋或者白雲把她們娶走，而對身邊千千萬普通男性充滿生理性厭惡，即使這些男性客觀條件比她們好得多。

而後者的典型，就是清華學姐誣陷案和追風小葉被誣陷案。女性在沒有強大男性庇護的時候，怕被普通男性侵犯的天性已經強大到讓她們極度焦慮、歇斯底里。

所以說實話，現代社會大齡剩女也是婦女解放和一夫一妻制的犧牲品，她們同時受無法找到優秀男性和害怕被普通男性騷擾的焦慮感折磨的痛苦，實在是太悲慘了。

在AI時代，AI軟件工程師無疑是全世界擁有最強生產力的一批人，也是戰鬥力最強的一批人。在此時此刻，我們漢族男性無疑已經有能力開發並量產超廉價AI自走殺戮機器人，可以在幾天內消滅一個國家的所有人口而不需要擔心核冬天（nuclear winter）的到來。假設美國能給願意移民美國的國男程式設計師一人發五個金髮白妹，美國能坐穩世界第一強國500年。可惜沒有一個國家這麼有眼光。

回到正題：雖然國男尤其是程式設計師和工人，是全世界真實實力最強的一批人，但是美國沒有給我們發五個金髮白妹，而國女也不會對真實的實力產生感覺。所以，如果你要讓人妻和未婚女性喜歡上我們，需要做的是在AI時代增強自身真實實力的同時，也要同時讓自己在女人眼中是一個原始社會中的強者。

接下來的一節，我會具體講述如何讓自己看着像一個原始社會中的強者，進而吸引母猴子。但是因爲一個人的時間精力是有限的，實際操作中，這兩者經常是矛盾的。至於多少時間用於增強自身真實實力，多少時間用於裝原始社會的強者來陪母猴子玩，交由讀者自行判斷。

\section{吸引力法則}

\subsection{永遠保持多偶}

\textbf{女人天性永遠無法理解爲什麼一個優秀的男人會保持單偶。}

一夫一妻制，在西歐僅僅實行了一千多年，而在包括中國在內的全世界所有其他地方都沒有根基。也就是說，女性根本沒有進化出一夫一妻對應的特質。所以在原始社會中，一個男人如果優秀，自然會吸引很多女人，那麼這個男人自然會選擇同時有幾個配偶。於是女人有一個天性，如果看到一個男人有很多或者很優秀的配偶（在現代就是女朋友或者老婆），這個人大概率是一個（在原始社會）優秀的男人，這個女性就很容易會不由自主地產生對這個男性的迷戀，想要也跟這個男性交配。女性的這種心理，有人稱之爲預選機制（preselection）。

當然，一個多偶的男人，和給女人提供安全感，是矛盾的。所以其實也是一個雙刃劍。之前我曾經做過大樣本的對照試驗，就是對一部分女人直接說我要多偶，對一部分女人不提想要多偶。結論其實很有趣，就是直接跟女人說要多偶的那個樣本，對我快速上頭的更多。尤其是我說我要多偶之後，她們往往會突然覺得我這個人很有趣，也很有種，然後就開始熱情聊天，快速上頭。但是負面是，當跟她們上牀之後，有些會安全感缺失，焦慮和嫉妒爆棚，然後主動跟我分手。然後每個人和我分手都是那種又愛又捨不得的痛苦不堪的感覺，看得我也好心疼。所以如果你想要比較長期的關係，這是一把雙刃劍，請小心使用。

而跟人妻交往就比較有意思了，因爲人妻有老公提供安全感，所以並不需要從黃毛身上追求安全感的感覺。所以有老婆有女朋友的黃毛，對於人妻只有附加的吸引力而沒有缺點。人妻可以盡情地和黃毛享受刺激的迷戀感，回去就可以享受老公提供的安全感。所以經常看到人妻的出軌對象也是已婚男人，也是這個原因。郭菊陽剛開始和我約會的時候，我同時有未婚的女朋友，我還把女朋友的照片給郭菊陽看，她看完更興奮了然後對我更加迷戀。沒有對未婚女性的那種副作用。

前面說到多偶有副作用，但是更有趣的是，單偶的副作用更大。如果你認定了一個女人，不再考慮和其他女人接觸的可能性，你在女人眼中就會完全喪失吸引力。女人會因爲缺乏安全感，一直跟你確認你能一心一意地爲她。比如會叫你官宣，叫你把你的手機屏保設成她的照片，不讓你接觸任何其他女性甚至和男性朋友出去約會。而一旦你真的滿足了她的這些所有的安全感需求，圍着她轉，她就會本能地覺得你竟然弱到能被她控制，不能按照自己本心去找多個女人，絕對是一個廢物。於是，她就會對你產生一個新的負面感覺：「無聊」。她會覺得你「沒刺激感」，「沒新鮮感」，對你「沒有感覺」。這時候，雖然你還是一心一意地對她好，她已經開始在外面追求新的基因提供者以獲得「感覺」了。

這就是爲什麼，你看很多又帥、又有才華、又有錢的完美丈夫，老婆還是會出軌（現代）或者離婚（老一輩）。在中國有汪小菲、焦恩俊，在國外有巴菲特、貝佐斯、比爾蓋茨，等等。無論男人多優秀，只要領證、保持安穩，女性就會沒有新鮮感和刺激感，想要找新的男人給她帶來新的刺激，即使客觀上這個男人比她優秀得多。而之後的出軌離婚，也便自然而然了。

有一次，我在深圳灣萬象城的至正潮菜宴請一個清華學弟，他當時博士畢業在騰訊做AI研發工作。他跟我說他之前好長一段時間都在接觸各種女生，終於最近交了個女朋友還挺喜歡的，就沒有繼續尋找了。當時我看着他的眼睛，很嚴肅地跟他說「\textbf{放棄多偶傾向，就是喪失男性魅力的開始。}」他突然安靜，沉思了很久，然後突然瞪大眼睛，對我露出了崇拜的眼神。大喊「太有哲理了」，差一點就認我做義父了。

而且，保持與多個女人接觸，事實上可以增加你能找到女朋友的期望總量。這也是很好理解的。假設每個女人$i$，你泡到她的概率是$p_i$。那麼你泡$1$到$n$一共$n$個女生，你期望得到的女朋友的數量就是$\sum_{i=1}^n p_i$。很明顯，你泡的女生越多，結果越好。如果你再把女人的價值$v_i$和泡她的成本$c_i$算進去，那麼你泡$n$個女人期望收益就是
$$\sum_{i=1}^n \left(p_i v_i - c_i\right)$$

這個建模還可以繼續細化，比如你的時間和金錢有限，所以你泡第一個女人和第十個女人的成本是不同的。所以你可以把成本分爲時間和金錢，然後就可以把泡到一個女人$i$的概率變成一個關於在這個女人身上花的時間$t_i$和金錢$w_i$函數$p_i(t_i, w_i)$，然後把你要泡那些女人變成一個有時間和金錢總量限制的凸優化（convex optimization）問題。

但是這不重要，這些公式其實都告訴我們，你看到一個女人，只要泡她的概率乘以收益，大於泡她的成本，你就應該直接去上去泡，而不需要被所謂「專一」的女權社會規訓所欺騙——這個社會女人從來就沒有專一過。現代社會，很多人覺得男人同時泡幾個女人就是「渣男」，缺德。而實際上，這不過是在做女人一直在做的事情罷了。很多喜歡把女人想象成聖女的「好男人」其實根本就不懂女人。女人看起來沒有主動出擊多個男人，但是女人只要通過發騷，比如在朋友圈或者小紅書發自拍，就可以等着一羣男人上鉤，然後再在主動聯繫的男人裏面挑那麼五個十個條件好、有「感覺」的保持曖昧關係。幾乎每個女人都會這麼做，而且她們還可以裝無辜說「我沒有主動是他們主動聯繫我的」，把自己的責任撇得一乾二淨。而女權分子很喜歡給人洗腦的「我可以騷你不能擾」，更是要靠輿論的力量強迫社會消除她們的任何責任。

所以男性廣撒網同時泡多個女性，無論是從自身利益出發還是從道德上的公平出發，都是毫無問題甚至是應該鼓勵的事情。

\subsection{永遠準備離開}

在和女人在一起的時候，要永遠準備隨時脫身。也就是說，永遠不要害怕失去一個女人。這點和多偶有點像，但是比多偶還要更間接、更抽象。就是你多偶的時候，一個女人讓你不開心，你就扭頭找另外一個女人去了。但是實際上，女人只要看到她對你不好你扭頭就走這一步，並不需要看到另外一個女人的存在，就會對你產生迷戀。一個優秀的男人，總是有自己的事情幹，不管是自己的興趣愛好、工作事業、還是別的女人。只有廢物的男人才會整天圍着女人轉。而一旦女人感受到你這種扭頭就走的魄力，就經常會產生迷戀。這種扭頭就走不理的行爲，有人稱之爲棄貓效應（abandoned cat effect）。

當年郭菊陽爲了嫁有錢人把我甩了，我就覺得她傻逼，於是幹自己的事情去了。結果到大二的時候我和正牌女友周遊全國的時候，她還跑去清華想找我。這就是棄貓效應的應用了。

\begin{figure}[H]
\centering
\includegraphics[width=3in]{figures/guo_7.jpg}
\caption{郭菊陽在高中被棄貓之後寫的}
\end{figure}

\begin{figure}[H]
\centering
\includegraphics[width=2.5in]{figures/guo_8.jpg}
\caption{棄貓效應持續到大二}
\end{figure}

最搞笑的是，很多女人似乎對自己會受棄貓效應有着清晰的認知，甚至以己度人，覺得對男人也有一樣的效果。我曾經認識過幾個女的，直接問我是不是她對我冷淡我就會喜歡她。我心裏笑着想，她就算直接撲上來我都可能會嫌棄，她要是冷淡的話，我大概率一輩子都不會理她了。

\subsection{家庭暴力}

首先聲明，無論是在中國還是西方國家，家庭暴力都是違法的。本章內容僅作爲學術探究，請讀者不要違反當地法律。

我有一個從海陸豐山區遷居到龍華山區的客家人小女朋友，每次見她，她都要求我用力地打她屁股。作爲一個從小到大無論到哪裏都遵紀守法的好（美國）公民，我第一次聽到的時候，非常爲難。我一開始還在想她是不是因爲土客械鬥失敗的歷史創傷造成了強迫性重複，作爲潮汕土人的我還有點愧疚。但是在她的強烈要求之下，還是用力地拍打她的屁股，結果她臉和屁股一樣變紅了，超級興奮。這時候我才知道打她是爲她好。讓樂於助人的我也變得特別開心和興奮。

我經常能看到新聞，說有丈夫每天毆打妻子，社工苦苦勸告妻子離婚，結果妻子看丈夫還是眼神充滿小星星，永遠不離不棄，甚至會罵社工不懷好意破壞家庭，然後小孩還是一個一個生。這種新聞不是個例，甚至有研究表明，在家暴之後，女性懷孕的概率顯著提高。

事實上女人天生會對有攻擊性的男生產生好感，即使她就是被攻擊的對象。有些女的會想，「他這麼有攻擊性，好有男人味，肯定能夠好好保護我。」「他爲什麼打我不是打別人，肯定心裏有我，把我當做他的女人，好有安全感。」

而其實相對於直接動手打女人，像我一樣罵女人也是一種比較間接的合法版本的家暴。比如郭菊陽，高中的時候她爲了嫁有錢人把我甩了開始，我就一直罵她賤貨了。到這幾年我每次操她的時候，也是打着屁股罵她是有錢就能上的賤貨。也算是被我罵爽十幾年了。

我也經常上網痛罵楊笠、痛罵張桂梅，很多人以爲女人們尤其是女權分子討厭死我了。其實我私下加過不少女權分子，她們很多面對我，都是那種戰戰兢兢，然後充滿尊敬地跟我說話。其實你看前面我說的「女權分子本質上是發泄她們慕強和缺乏安全感的情緒」就懂了。當她們被一個強者指着鼻子痛罵，她們也就被罵爽了，罵舒服了，不想打拳了。要是我願意操她們，明天就開始去商場挑小孩子衣服了，還打什麼拳？

三年前有一個舔狗衆多，父親是公安局長的女羣友。我約她去KTV，結果她一副跟我去玩是給我恩賜的表情，最後還爽約了。因爲我曾經被網警拘留過，本身就看她這種公安二代很不爽。結果她還這麼不尊重我，於是我就在羣裏罵了她兩年。結果去年有一次，陰差陽錯友好地請她喫了一次飯，她就瘋狂地迷戀上我了。讓我非常莫名其妙。或許這就是家庭（語言）暴力的魅力了。

\subsection{革命家和藝術家}

其實比起家暴她的女人，有一種女人更喜歡的有攻擊性的性格，我稱之爲「革命家性格」。革命家不僅僅要有攻擊性，還要有正義感和發自內心的勇氣，敢於對抗所有的強權。而最看不起的，就是那些爲了錢失去原則，什麼人都能跪，什麼生意都能做的小人。

女人會瘋狂迷戀的另一種人，就是藝術家。像音樂、舞蹈，本身就是人類最原始的情感表達，比什麼都直達人心。而美術與文學，更是能創造出不朽的作品，對抗時間的虛無，把人類情感凝聚成永恆。又有什麼比藝術創作更偉大而值得獻身的呢？而我，不也不吝耗費自己本來能用來賺幾千萬幾億的寶貴時間，凝聚成《人妻約會指南》這本藝術作品嗎？

歷史上的革命家和藝術家，從來不缺女人不要求金錢不要求名分，飛蛾撲火一般撲上去獻身。即使在2025年的今天，走在紐約的街頭，還有無數的白妹穿着切格瓦拉的T恤。唱出《怒放的生命》的汪峯，有着三個美貌的老婆和4個兒女。而唱《一無所有》的崔健，小女友也沒間斷過，幾十年前就有女友未婚爲他生下一個漂亮的女兒。

有時候我在想，革命家、藝術家，對個體來說，經常都是風險巨大且經常會窮困潦倒的職業。但是一個族羣，如果沒有革命家，沒有藝術家，那麼這個族羣，還有什麼偉大的可能，又有什麼存在的意義？女性對革命家和藝術家的飛蛾撲火，何嘗又不是一種族羣至上的蜂巢思維呢？

我的父親李松堅，作爲一個低聲下氣跪着被雙開的陳良宇、蔣超良賺黑錢的黑心裙帶資本家（crony capitalist），被凌菲菲戴綠帽、背刺、霸佔財產，真的是太合理了。又有哪個女人會真心喜歡這種醜陋猥瑣的革命家的對立面呢？而凌菲菲拿到幾十億之後，也去追求她的真愛，原國家博物館館長，畫家陳履生去了。

\subsection{保持主體性}

其實上面分析了這麼多，其實吸引力法則總結起來就是一句話：\textbf{忘記女人，做你自己。}

錢總是流向不缺錢的人，而女人總是撲向那些不缺女人的人。你如果以想要女人甚至是想要操逼的心態來讀上面的章節，你一定是學不像的，最後總是會透露出猥瑣的氣息。

而要獲得吸引力，就是要保持主體性，把注意力放在自己身上：健身、學習、創作、改變世界。誰敢侵犯你，你就破口大罵。哪個女朋友讓你不舒服了，你扭頭就走。把人生去追求科學和藝術等永恆的東西。如果你發自內心地愛自己、做自己、追求自己的理想，那麼你自然而然能有最強大的吸引力了。\footnote{可以參考後面的章目《超越人妻》。}

\section{如何吸引人妻}

人妻和單身的女人相比，其實更好吸引。首先，對於一個現代的單身女性，她們既要你能夠給她們迷戀的感覺，也就是需要讓她們覺得你是強者，又要你給她們一定的安全感，也就是需要你有錢適當供養她們。而對人妻來說，絕大部分對安全感的需求都可以由她們老公來滿足，所以你只要讓她們產生迷戀的感覺就行了。

其實很多時候，讓女人產生迷戀的感覺並不難，其實比長期維持關係要簡單。因爲一旦結婚長期生活在一起，丈夫的很多優點，就被當作理所當然，而丈夫的缺點，就會被一直注意到。這時候，就會造成一種強烈的缺失感和需要補償的感覺。郭菊陽當着我的面說過，她後悔大學的時候沒有多跟幾個男的上牀，感覺喫虧了。類似的話，我也不止一次聽不同的女人說過。每個結婚的女人都覺得虧了，因爲沒法多試試幾個男人。所以你只要多聊幾個人妻，如果你剛好滿足了某個人妻在當下的缺失感，即使你整體遠遠不如苦主，人妻也會想跟你談戀愛上牀。

話說人妻們對老公挑刺，是可以極其離譜的。她們罵老公的話，你們一句話都別信。我在美國讀研究生的時候每天陪我喫飯的女飯搭子，後來嫁到宇宙中心加州帕羅奧圖（Palo Alto）。她整天給我發短訊、微信、WhatsApp，說她老公不賺錢她生活很苦。後面我跟他老公聊天，才知道「老公不賺錢」是稅前賺二三十萬美金現金，而她一年要花稅後五十萬美金的意思。知道之後，我非常嚴肅地替她老公批評教育了她。結果只換來了一句「原來你也是惡臭摳門男」的評價。

但是讀者們站在一個黃毛的角度，就知道這些有錢人的嬌妻們是多麼容易挑刺和不滿，我們能找到的漏洞何其多！如果苦主年齡大，那麼人妻就會想要找年輕的出軌；如果苦主長得醜，那麼人妻就會找長得帥的出軌；如果苦主賺錢少，那麼人妻就會找高收入黃毛出軌；如果苦主學歷低，那麼人妻就會找高學歷的出軌；如果苦主在外面包二奶，那麼人妻就會隨便出軌；如果人妻就是苦主包的二奶，那也會隨便出軌！

所以即使你是一個3分男人，你就同時撩那麼100個7分男人的老婆，展示自己的優點。總會有時候，你的某個優點，剛好是某個人妻心理缺失的那一塊，這時候她就會對你上頭、迷戀，把你當做基因提供者了。

說實話，我泡郭菊陽的經歷沒什麼好講的，因爲我沒什麼短板，說自己經歷有點「數值怪以爲自己可有操作了」的感覺。你想泡李松堅領證的好老婆許潔\footnote{見附錄\hyperref[sister]{《我的姐姐李新瑩》}}的話，也是沒什麼好講的，畢竟李松堅樓盤爛尾被銀行抽貸。一個長得像曹德旺的負資產的60多歲的中學學歷的老頭，操他老婆又有什麼難的呢？

他老婆如果大家喜歡的話，大家閉着眼睛直接衝就是了。如果不喜歡的話，想到她應該靠長期做雞最後還傍上大款上位成功賺了那麼多錢，大家看在錢的份上也努力衝吧。許潔的電話我也不知道哪個，大家自己試試：13585833555、18019099644、18019243004、18101963394。

\begin{figure}[H]
\centering
\includegraphics[width=1.2in]{figures/xu.jpg}
\caption{「教授雞」許潔在上海音樂學院官網的照片}
\end{figure}
